\section{Conclusion}
\label{sec:conclusion}

We presented \method{}, a pipeline that turns recorded agent-execution failures into reusable defenses: trajectory graphs feed a label-free anomaly detector, a weakly supervised node scorer localizes the compromised span, an LLM judge distills the localized evidence into structured safety experiences under a strict no-label, redaction-first evidence boundary, and a paired replay protocol measures the resulting protection as a within-case ASR reduction with an independent judge family. On \tool{}, the pipeline detects compromised trajectories across five seeds, converts detector-flagged cases into admitted experiences at a principled rate, and reduces attack success relative to the paired baseline, with ablations attributing the improvement to graph-derived localization rather than to generic safety instructions. The approach occupies a defense point---post-execution memory injection---that complements prompt hardening, information-flow separation, and runtime monitoring, and it inherits an explicit coverage/recall trade-off from its detector front end rather than hiding it. Future work includes cross-dataset and cross-model experience transfer, human-audited judge calibration at scale, adaptive-attack robustness of the replay protocol, and integrating the experience store with the persistent memory systems now emerging for long-horizon agents.
